% Zdroj: PDF priklady-leto-2022-one.pdf (sešit S1, PGVVH+UI), fyzická str. 1-2. Značka: společné okruhy. Zařazeno: I3 (překlad a běh). Číselné/logické výsledky ověřeny.
\section{Překlad konstrukcí vyššího jazyka}
\szzotazka{léto 2022}{2}{Překlad konstrukcí vyššího jazyka}

\noindent\textbf{Zadání.}\par
Mějme jednoduchý procesor s architekturou inspirovanou MIPS. Procesor má 32
32-bitových celočíselných obecných registrů \texttt{\$R0}--\texttt{\$R31} a jeden
32-bitový registr \texttt{PC} ukazující na začátek následující instrukce. Hodnota
registru \texttt{\$R0} je vždy 0. V popisu operandů značí \texttt{label} symbolické
návěští, \texttt{[mem]} adresu paměti a \texttt{reg} obecný registr.

\medskip
\noindent Instrukce procesoru:
\begin{itemize}
  \item \texttt{LD reg=[mem]} -- načte celočíselnou proměnnou z adresy \texttt{mem}
    do registru (\texttt{reg = *mem}).
  \item \texttt{ST [mem]=reg} -- uloží registr do proměnné na adrese \texttt{mem}
    (\texttt{*mem = reg}).
  \item \texttt{ADD reg1=reg2,reg3} -- \texttt{reg1 = reg2 + reg3}.
  \item \texttt{J label} -- nepodmíněný skok (\texttt{PC = label}).
  \item \texttt{BEQ reg1,reg2,label} -- skok, pokud \texttt{reg1 == reg2}.
  \item \texttt{BNE reg1,reg2,label} -- skok, pokud \texttt{reg1 != reg2}.
  \item \texttt{SLT reg1=reg2,reg3} -- \texttt{reg1 = (reg2 < reg3)} (jinak 0).
\end{itemize}

\medskip
\noindent Mějme následující posloupnost instrukcí:
\begin{lstlisting}
        LD    $r1=[q]
        LD    $r2=[n]
        LD    $r3=[inc]
        J     lab1
lab2:
        LD    $r4=[a]
        LD    $r5=[b]
        ADD   $r4=$r4,$r1
        SLT   $r5=$r5,$r4
        BEQ   $r5,$r0,lab3
        ST    [b]=$r4
lab3:
        ADD   $r1=$r1,$r3
lab1:
        SLT   $r4=$r1,$r2
        BNE   $r4,$r0,lab2
\end{lstlisting}

\begin{enumerate}
\item Vyberte ze zadaných úryvků v jazyce C všechny, které odpovídají výše uvedené
  posloupnosti instrukcí (neznámé identifikátory považujte za někde deklarované
  32-bitové celočíselné proměnné). Pro zvolené řešení připište k~instrukcím čísla řádků
  odpovídajících příkazů v~jazyce C. Nabízené úryvky (v~PDF je každý rozepsán na očíslované
  řádky, jeden příkaz nebo hlavička na řádek: úryvky (a)--(d) mají řádky 1--3, úryvek (e) řádky 1--6):

  \begin{enumerate}
    \item \texttt{for (i=q; i<n; i+=inc) if (a+i <= b) b = a+i;}
    \item \texttt{if (a+i > b) for (i=q; i<n; i+=inc) b = a+i;}
    \item \texttt{for (i=q; i<n; i+=inc) if (a+i > b) b = a+i;}
    \item \texttt{if (a+i < b) for (i=q; i<n; i+=inc) b = a+i;}
    \item \texttt{i = q; while (i<n) \{ if (a+i > b) b = a+i; i += inc; \}}
    \item Žádný z uvedených úseků neodpovídá (napište vlastní kód).
  \end{enumerate}

\item Napište (stačí jeden řádek), jak se původní kód v jazyce C změní, pokud
  instrukci \texttt{BEQ} nahradíme instrukcí \texttt{BNE} se stejnými operandy.

\item Předpokládejme, že podmínka v příkazu \texttt{if} je rozšířena o část
  \texttt{\&\& i < b}, tj. bude mít tvar \texttt{if (původní podmínka \&\& i < b)}.
  Upravte původní sekvenci instrukcí tak, aby tuto rozšířenou podmínku realizovala.
\end{enumerate}

\medskip
\noindent\textbf{Náčrt řešení.}\par
\textit{(V~PDF bez náčrtu řešení; náčrt doplněn k~výkladu okruhu, logika ověřena
simulací.)}

\medskip
Nejprve rozluštíme posloupnost. Před cyklem se načte \texttt{q}, \texttt{n},
\texttt{inc} (do \texttt{\$r1}, \texttt{\$r2}, \texttt{\$r3}) a skočí se na
\texttt{lab1}. Tam je test cyklu: \texttt{SLT \$r4=\$r1,\$r2} nastaví
\texttt{\$r4 = (i < n)} (řídicí proměnná \texttt{i} je v \texttt{\$r1}, počáteční
hodnota \texttt{q}) a \texttt{BNE \$r4,\$r0,lab2} skočí do těla, je-li
\texttt{\$r4 != 0}, tj. \emph{dokud platí} \texttt{i < n}. V těle se spočítá
\texttt{\$r4 = a + i}, pak \texttt{SLT \$r5=\$r5,\$r4} dá \texttt{\$r5 = (b < a+i)}
a \texttt{BEQ \$r5,\$r0,lab3} \emph{přeskočí} uložení \texttt{ST [b]=\$r4}, když
\texttt{\$r5 == 0}, tj. když \texttt{a+i <= b}. Uložení \texttt{b = a+i} se tedy
provede právě tehdy, když \texttt{a+i > b}. Na \texttt{lab3} se \texttt{i} zvýší
o \texttt{inc} a běh pokračuje testem cyklu.

\medskip
\noindent\textbf{(1)} Posloupnost realizuje
\begin{lstlisting}[language=C]
for (i=q; i<n; i+=inc)
    if (a+i > b) b = a+i;
\end{lstlisting}
a~ekvivalentně její while-rozpis
\begin{lstlisting}[language=C]
i = q;
while (i<n) { if (a+i > b) b = a+i; i += inc; }
\end{lstlisting}
Odpovídají tedy \emph{obě} varianty: \emph{(c)} for s podmínkou \texttt{a+i > b}
i~\emph{(e)} ekvivalentní while. Varianty (a) (opačná podmínka \texttt{<=}),
(b)/(d) (zaměněné pořadí for a~if, resp. ostré \texttt{<}) neodpovídají.

\medskip
\noindent Čísla řádků k~instrukcím. Úryvek (c): řádek~1 \texttt{for (i=q; i<n; i+=inc)},
2~\texttt{if (a+i > b)}, 3~\texttt{b = a+i;}. Úryvek (e): 1~\texttt{i = q;}, 2~\texttt{while (i<n) \{},
3~\texttt{if (a+i > b)}, 4~\texttt{b = a+i;}, 5~\texttt{i += inc;}, 6~\texttt{\}}.
\begin{center}
\begin{tabular}{lcc}
instrukce & (c) & (e) \\ \hline
\texttt{LD \$r1=[q]} & 1 & 1 \\
\texttt{LD \$r2=[n]}, \texttt{LD \$r3=[inc]} & 1 & 2, 5 \\
\texttt{J lab1} & 1 & 2 \\
\texttt{LD \$r4=[a]}, \texttt{LD \$r5=[b]}, \texttt{ADD \$r4=\$r4,\$r1}, \texttt{SLT \$r5=\$r5,\$r4}, \texttt{BEQ} & 2 & 3 \\
\texttt{ST [b]=\$r4} & 3 & 4 \\
\texttt{ADD \$r1=\$r1,\$r3} & 1 & 5 \\
\texttt{SLT \$r4=\$r1,\$r2}, \texttt{BNE} & 1 & 2 \\
\end{tabular}
\end{center}
Součet \texttt{a+i} v~\texttt{\$r4} spočítaný pro podmínku se znovu použije jako přiřazovaná
hodnota na řádku s~\texttt{b = a+i} (proto je \texttt{ADD \$r4=\$r4,\$r1} vlastně společný
podmínce i~přiřazení). Test \texttt{i<n} je v~posloupnosti \emph{na konci} cyklu a~\texttt{J lab1}
na něj skočí hned na začátku, což je obvyklý překlad cyklu \texttt{for}/\texttt{while}.

\medskip
\noindent\textbf{(2)} Záměna \texttt{BEQ} za \texttt{BNE} (se stejnými operandy
\texttt{\$r5,\$r0,lab3}) obrátí podmínku skoku: přeskočení uložení nastane, když
\texttt{\$r5 != 0}, tj. když \texttt{a+i > b}. Uložení se proto provede v opačném
případě, takže podmínka v~C se změní na
\begin{lstlisting}[language=C]
if (a+i <= b) b = a+i;
\end{lstlisting}

\medskip
\noindent\textbf{(3)} Rozšíření o \texttt{\&\& i < b}: před instrukci
\texttt{ST [b]=\$r4} vložíme druhý test. Načteme \texttt{b} (do \texttt{\$r6}),
spočítáme \texttt{\$r6 = (i < b)} a \texttt{BEQ \$r6,\$r0,lab3} přeskočí uložení,
když \texttt{i < b} neplatí. Uloží se tedy jen při splnění obou podmínek
(\texttt{a+i > b} a zároveň \texttt{i < b}). Upravená posloupnost (komentáře bez
diakritiky):
\begin{lstlisting}
lab2:
        LD    $r4=[a]
        LD    $r5=[b]
        ADD   $r4=$r4,$r1
        SLT   $r5=$r5,$r4        ; r5 = (b < a+i)
        BEQ   $r5,$r0,lab3       ; preskoc, kdyz a+i <= b
        LD    $r6=[b]            ; pridany test i < b
        SLT   $r6=$r1,$r6        ; r6 = (i < b)
        BEQ   $r6,$r0,lab3       ; preskoc, kdyz i < b neplati
        ST    [b]=$r4            ; ulozi se jen kdyz a+i > b AND i < b
lab3:
        ADD   $r1=$r1,$r3
\end{lstlisting}
